\documentclass[10pt]{article}
\providecommand{\FIGDIR}{fig}
\newcommand{\DIGESTDATE}{29 August 2026}
\newcommand{\PERIODSTART}{23 August 2026}
\newcommand{\PERIODEND}{29 August 2026}
% ===================== PM Research Watch - shared preamble =====================
\usepackage[T1]{fontenc}
\usepackage[utf8]{inputenc}
\usepackage{amsmath,amssymb}
\usepackage{tgpagella}
\usepackage{tgheros}
\renewcommand{\sfdefault}{qhv}
\renewcommand{\ttdefault}{lmtt}
\usepackage{microtype}
\usepackage[english]{babel}

\usepackage[a4paper,top=17mm,bottom=17mm,left=17mm,right=17mm,
            headsep=6mm,footskip=10mm]{geometry}
\usepackage{graphicx}
\usepackage{xcolor}
\usepackage{colortbl}
\usepackage{tikz}
\usetikzlibrary{calc}
\usepackage[most]{tcolorbox}
\usepackage{titlesec}
\usepackage{fancyhdr}
\usepackage{booktabs}
\usepackage{array}
\usepackage{longtable}
\usepackage{enumitem}
\usepackage{multicol}
\usepackage{ragged2e}
\usepackage[hidelinks]{hyperref}

% ------------------------------------------------------------------ palette
\definecolor{Ink}    {HTML}{14262E}
\definecolor{Deep}   {HTML}{0F4C5C}
\definecolor{Teal}   {HTML}{1F7A8C}
\definecolor{Sky}    {HTML}{4FA3C4}
\definecolor{Sage}   {HTML}{6FA287}
\definecolor{Amber}  {HTML}{E3A008}
\definecolor{Clay}   {HTML}{B5643C}
\definecolor{Coral}  {HTML}{D1495B}
\definecolor{Violet} {HTML}{6A5B9E}
\definecolor{Slate}  {HTML}{4A6572}
\definecolor{Paper}  {HTML}{FBF9F5}
\definecolor{Mist}   {HTML}{EEEAE2}
\definecolor{Rule}   {HTML}{D9D3C9}

\pagecolor{Paper}
\color{Ink}

\hypersetup{colorlinks=true, linkcolor=Deep, urlcolor=Teal, citecolor=Teal,
            pdftitle={PM Research Watch}, pdfauthor={Automated daily digest}}

% ------------------------------------------------------------------ headings
\titleformat{\section}
  {\sffamily\bfseries\large\color{Deep}}{}{0pt}
  {}[{\vspace{1.2mm}\color{Amber}\titlerule[1.1pt]}]
\titlespacing*{\section}{0pt}{5.5mm}{2.8mm}

\titleformat{\subsection}
  {\sffamily\bfseries\normalsize\color{Teal}}{}{0pt}{}
\titlespacing*{\subsection}{0pt}{3.6mm}{1.4mm}

\setlist[itemize]{leftmargin=4.2mm, itemsep=1.1mm, topsep=1.3mm, parsep=0pt,
                  label=\textcolor{Teal}{\footnotesize$\blacktriangleright$}}

% ------------------------------------------------------------------ header/footer
\pagestyle{fancy}
\fancyhf{}
\renewcommand{\headrulewidth}{0pt}
\renewcommand{\footrulewidth}{0pt}
\fancyhead[L]{\sffamily\scriptsize\color{Slate}PM RESEARCH WATCH\;\textcolor{Rule}{\textbar}\;\DIGESTDATE}
\fancyhead[R]{\sffamily\scriptsize\color{Slate}Particulate matter monitoring \& health effects}
\fancyfoot[C]{\sffamily\scriptsize\color{Slate}\thepage}

% ------------------------------------------------------------------ components
% metric tile
\newcommand{\metric}[3]{%
  \begin{tcolorbox}[enhanced, colback=white, colframe=Rule, boxrule=0.5pt,
    arc=1.6mm, left=2.2mm, right=2.2mm, top=1.8mm, bottom=1.8mm,
    borderline west={1.7mm}{0pt}{#3}, nobeforeafter, valign=center,
    height=17mm, width=\linewidth]
    \RaggedRight
    {\sffamily\bfseries\fontsize{17}{17}\selectfont\color{#3}#1}\\[0.5mm]
    {\sffamily\fontsize{6.5}{7.4}\selectfont\color{Slate}#2}
  \end{tcolorbox}}

% pull-quote / key insight box
\newtcolorbox{keybox}[1][]{enhanced, colback=Mist, colframe=Mist,
  arc=1.6mm, left=3.2mm, right=3.2mm, top=2.4mm, bottom=2.4mm,
  borderline west={1.7mm}{0pt}{Amber}, #1}

% subtopic band
\newcommand{\band}[2]{%
  \vspace{3.2mm}\noindent
  \begin{tcolorbox}[enhanced, colback=#2, colframe=#2, arc=1.2mm,
    left=2.6mm, right=2.6mm, top=1.1mm, bottom=1.1mm, boxrule=0pt,
    width=\linewidth, nobeforeafter]
    {\sffamily\bfseries\small\color{white}#1}
  \end{tcolorbox}\vspace{1.6mm}}

% one paper entry: {short cite}{title}{journal}{doi}{body}
\newcommand{\paperentry}[5]{%
  \par\vspace{1.9mm}\noindent
  {\sffamily\bfseries\footnotesize\color{Ink}#1}\;%
  {\sffamily\fontsize{7}{8}\selectfont\color{Slate}#3}\\[0.4mm]
  {\footnotesize\itshape\color{Deep}#2}\\[0.6mm]
  {\footnotesize #4}\\[0.3mm]
  {\fontsize{6.6}{7.6}\selectfont\color{Slate}\texttt{doi:}\,\href{https://doi.org/#5}{\texttt{\detokenize{#5}}}}
  \par}

% caption style
\newcommand{\figcap}[1]{\par\vspace{1.0mm}
  {\fontsize{7.2}{8.6}\selectfont\color{Slate}\textsf{\textbf{Fig.}}\; #1\par}\vspace{1.5mm}}

\providecommand{\permil}{\textperthousand}

\setlength{\parindent}{0pt}
\setlength{\parskip}{0pt}
\renewcommand{\arraystretch}{1.18}


\begin{document}

% ==================================================================== MASTHEAD
\thispagestyle{empty}
\begin{tikzpicture}[remember picture, overlay]
  \fill[Deep] (current page.north west) rectangle
      ($(current page.north east)+(0,-67mm)$);
  \fill[Teal] ($(current page.north west)+(0,-67mm)$) rectangle
      ($(current page.north east)+(0,-68.4mm)$);
  \foreach \x/\y/\r in {22/12/0.9, 41/26/1.5, 63/9/0.7, 88/31/1.1, 112/17/1.9,
                        138/29/0.8, 159/13/1.3, 178/54/1.0, 196/21/1.6,
                        30/44/0.6, 74/57/1.2, 125/48/0.9, 168/6/0.8, 191/59/0.7,
                        52/18/1.0, 100/60/0.6, 148/20/1.1, 66/49/1.4, 205/45/0.8}
    \fill[Sky, opacity=0.30] ($(current page.north west)+(\x mm,-\y mm)$) circle (\r mm);
\end{tikzpicture}

\vspace*{4mm}
{\sffamily\fontsize{7.6}{9}\selectfont\color{Sky}%
 WEEKLY ROLLUP\;\textbullet\;PhD RESEARCH WATCH\par}
\vspace{2.2mm}
{\sffamily\bfseries\fontsize{27}{29}\selectfont\color{white}%
 Particulate Matter\par}
{\sffamily\fontsize{27}{29}\selectfont\color{Sky}%
 Week in Review\par}
\vspace{5.0mm}
{\sffamily\fontsize{8.4}{10}\selectfont\color{white}%
 \mbox{Week ending \DIGESTDATE}\;\;\textcolor{Teal}{\textbar}\;\;\mbox{\PERIODSTART{} to \PERIODEND}\;\;
 \textcolor{Teal}{\textbar}\;\;\mbox{109 records pooled from 7 daily issues}\par}
\vspace{18mm}

% -------------------------------------------------------------- metric strip
\noindent
\begin{minipage}[t]{0.190\linewidth}\metric{109}{RECORDS POOLED\\(107 UNIQUE DOIs)}{Deep}\end{minipage}\hfill
\begin{minipage}[t]{0.190\linewidth}\metric{41\%}{ON THE MEASUREMENT\\SIDE (45 RECORDS)}{Teal}\end{minipage}\hfill
\begin{minipage}[t]{0.190\linewidth}\metric{10}{SUBTOPICS\\ACTIVE}{Sky}\end{minipage}\hfill
\begin{minipage}[t]{0.190\linewidth}\metric{37}{POOLED EFFECT\\ESTIMATES}{Amber}\end{minipage}\hfill
\begin{minipage}[t]{0.190\linewidth}\metric{8}{ABSTRACT-FREE\\(7\% OF WEEK)}{Clay}\end{minipage}

\vspace{6mm}

% -------------------------------------------------------------- week in one box
\section*{\sffamily\bfseries\color{Deep}The week in one paragraph}
\vspace{-1.5mm}{\color{Teal}\rule{\linewidth}{1.1pt}}\vspace{2.5mm}

\begin{keybox}
{\sffamily\bfseries\small\color{Deep}The health side took the week back from the
measurement side, and the abstract-free share collapsed from 20\% to 7\% --- but the
week's real content is a convergence: three independent lines all say the health-relevant
quantity is a \emph{composition}, and none of them is measurable by the instruments the
epidemiology relies on.}\\[1.4mm]
\footnotesize
\textbf{The composition thread runs the week end to end.} Li et al.\ (28 Aug,
\textit{Environ Sci Technol Lett}) decompose US county AMI mortality into PM$_{2.5}$
constituents and find \textbf{SO$_4^{2-}$ carrying weight 0.762} of a mixture effect
(RR 1.116 per quartile) that is several times the total-mass estimate (RR 1.019 per IQR),
with the dominant constituent flipping east-to-west from sulfate to organic carbon. Wang
et al.\ (28 Aug, \textit{Sleep Med}) run the same five-constituent decomposition through
19{,}505 CHARLS participants and land on the secondary inorganic ions again. Zhang et al.\
(28 Aug, \textit{ACP}) supply the physics: single-particle SP2 and SPAMS measurement in
Shenzhen shows nighttime heterogeneous and daytime photochemical aging arrive at nearly
the same \textbf{mass absorption cross-section enhancement --- $\sim$0.8 and
$\sim$1.0\,m$^2$\,g$^{-1}$ at 532\,nm} --- by chemically different routes. A mass-only
optical monitor re-expresses that diurnal coating cycle as concentration drift.

\textbf{The counter-observation is that the coarse fraction refuses to go away.}
Park et al.\ (29 Aug) put \textbf{PM$_{10}$ ahead of PM$_{2.5}$} for the
high-TSH/low-FT4 thyroid phenotype (aOR 1.55 vs.\ 1.25 per IQR), and Cliff et al.\ (29 Aug)
identify atmospheric mineral dust as the strongest single predictor of meningococcal
outbreaks across the African meningitis belt. Neither is causally identified --- both rest
on collinear exposure fields --- but the fine-fraction default is doing less work in this
week's data than the literature's framing assumes.

\textbf{Operationally the week improved.} \textbf{Abstract-free deposits fell from 22 of
112 (20\%) to 8 of 109 (7\%)}, so \textit{Metadata only} is no longer the largest
architecture label; \textit{Measurement campaign} (24) is. Two defects offset that: two
records were published twice under different daily issues, and five short-label
collisions were created and repaired inside this period. Both are recorded below.
\end{keybox}

\clearpage

% -------------------------------------------------------------- trend
\section*{\sffamily\bfseries\color{Deep}Subtopic trend and inflection}
\vspace{-1.5mm}{\color{Teal}\rule{\linewidth}{1.1pt}}\vspace{2.5mm}

\noindent\includegraphics[width=\linewidth]{\FIGDIR/w1_subtopic_trend.png}
\figcap{Per-issue subtopic counts across the week from \texttt{state/metrics.csv}.
\textbf{Read the denominators first}: daily totals were 9, 9, 14, 14, \textbf{36}, 19 and
8 --- a 4.5$\times$ range, far more volatile than last week's 2.2$\times$. The 27 August
figure of 36 is the largest single-day issue this watch has produced outside the founding
backfill, and it is a publisher artefact: a single MDPI \textit{Toxics} volume release
contributed six records into a 43-deposit Crossref day. \textbf{The trend line for
27 August therefore cannot be read at face value}, and every week-over-week movement below
should be discounted accordingly.}

\footnotesize
\textbf{The measurement side gave back last week's gain.} \textit{Exposure assessment \&
modelling} fell \textbf{36 $\rightarrow$ 24} while \textit{Sensing, forecasting \&
instrumentation} held at \textbf{22 $\rightarrow$ 21}; together they are 45 of 109, a
measurement-side share of \textbf{41\%, against 52\%, 45\% and 54\% in the three preceding
weeks}. The four-week sequence 54, 45, 52, 41 has no trend in it --- it is the range this
pipeline's own harvesting legs produce as the Crossref by-ISSN sweep competes with PubMed
for share of a given day. No claim about the field should rest on it.

\vspace{1.2mm}
\textbf{Expanding.} \textit{Burden, policy \& mitigation} rose \textbf{9 $\rightarrow$ 13}
and \textit{Mechanistic toxicology} \textbf{8 $\rightarrow$ 12}, the latter reversing last
week's break in a two-week run at 11. \textit{Neuro / mental health}
(\textbf{4 $\rightarrow$ 7}), \textit{Occupational \& indoor}
(\textbf{4 $\rightarrow$ 7}) and \textit{Cardiovascular \& metabolic}
(\textbf{3 $\rightarrow$ 5}) all recovered from lows that the 22 August rollup flagged as
worth a second week of watching; cardiovascular at 5 is still below its running level and
still well under the respiratory count.

\vspace{1.2mm}
\textbf{Contracting.} \textit{Other clinical endpoints} fell \textbf{13 $\rightarrow$ 8},
partially unwinding the two-week finding that the clinical tail had outgrown the
cardiorespiratory trunk. \textit{Reproductive \& developmental} at 3 is the week's
smallest cluster and its lowest reading this month --- notable against last week, when the
developmental mixture cohorts supplied most of the constituent-resolved estimates.
\textbf{All ten subtopics were active on at least one day, for the fifth consecutive
week.}

%% LAYOUT RULE (2026-08-02): architecture donut in the whitespace here, not on its own page.
\vspace{3mm}
\noindent\includegraphics[width=\linewidth]{\FIGDIR/f2a_architecture.png}
\figcap{Pooled study architecture over 109 records. \textbf{\textit{Metadata only} has
fallen from 22 records (20\% of last week) to 8 (7\%)} and is no longer the largest
label --- the title-pattern front-matter filter added on 22 August, plus a deliberate
decision across this week's issues to reject abstract-free deposits to a retry list rather
than carry them at tier C, both contributed. \textit{Measurement campaign} (24) now leads,
ahead of \textit{Review / synthesis} (17) and \textit{Modelling / inventory} (14);
\textit{Observational --- cohort} (11) and \textit{Experimental / toxicology} (9) are the
load-bearing evidence classes. \textbf{\textit{Trial / intervention} is 0 records in
109} --- for the fifth week running no completed randomised PM intervention entered the
corpus, and this week not even a process evaluation did.}

\clearpage

% -------------------------------------------------------------- pooled analytics
\section*{\sffamily\bfseries\color{Deep}Pooled analytics}
\vspace{-1.5mm}{\color{Teal}\rule{\linewidth}{1.1pt}}\vspace{3mm}

\noindent\includegraphics[width=\linewidth]{\FIGDIR/f1_subtopics.png}
\figcap{Pooled subtopic distribution, 23--29 August. \textit{Exposure assessment \&
modelling} (24) and \textit{Sensing} (21) are 41\% of the week between them ---
the lowest measurement-side share in four weeks. \textit{Burden, policy \& mitigation}
(13) is the largest health-adjacent cluster, ahead of \textit{Mechanistic toxicology}
(12) and \textit{Respiratory \& allergic} (9). Every one of the ten canonical clusters is
occupied, now for five consecutive weeks.}

\vspace{2mm}
\noindent\includegraphics[width=\linewidth]{\FIGDIR/f2b_endpoint.png}
\figcap{Health endpoint by organ system across the pooled week. \textbf{39 of 109 records
carry no human health endpoint} --- roughly 8 because no abstract was deposited and the
endpoint could not be extracted, the rest because they are monitoring, exposure,
instrument or aerosol-chemistry work by design. That is 36\% against last week's 55\%, and
the shift is the direct consequence of the abstract-free collapse plus the health-side
rebound in the cluster counts. \textbf{Endpoint labels remain uncanonicalised}, so the
long tail of singleton endpoint strings in this panel overstates endpoint diversity ---
see the disclosed defect on the venues page.}

\vspace{2mm}
\noindent\includegraphics[width=\linewidth]{\FIGDIR/f3_metric_geography.png}
\figcap{Left --- particle metric across the pooled week. PM$_{2.5}$ alone is the exposure
in 37 records, unspeciated PM or emissions in 36, and PM$_{2.5}$ with PM$_{10}$ jointly in
23. \textbf{Ultrafine or nanoscale particles are the explicit subject of 4 records and
size-resolved number distributions of 3 more} --- 7 of 109, double last week's 2 of 112
and the highest this watch has recorded, though still an order of magnitude below what
deposition physics would justify. Only 2 records carry a \textit{composition / speciation}
label, which is the standing grouper limitation: the week's constituent-resolved
epidemiology bins as \textit{PM$_{2.5}$ only}. Right --- geographic provenance:
\textbf{multi-region syntheses (28) edged past China (27)} for the first time in this
watch, with Europe 19 and North America 10. Sub-Saharan Africa contributes 2, both on
29 August, after an eight-issue absence.}

\clearpage

\noindent\includegraphics[width=\linewidth]{\FIGDIR/f5_forest.png}

\vspace{1mm}
\noindent\includegraphics[width=\linewidth]{\FIGDIR/f5_forest_2.png}
\figcap{All 37 pooled effect estimates, split across two panels and plotted on a common
logarithmic axis. \textbf{Exposure contrasts are heterogeneous and the estimates are not
mutually comparable} --- per-IQR, per-10\,$\mu$g/m$^3$, per-SD, quartile contrasts and
binary residential exposures all appear --- so read direction and precision, never
relative magnitude, and \textbf{no meta-analytic summary is computed or implied}. Two
internally comparable sub-sets exist and are worth separating from the rest. The five
Park et al.\ thyroid estimates (29 Aug) share one outcome, one cohort and one per-IQR
contrast, so their ordering is internally consistent --- but they come from a single
collinear CMAQ field entered in separate single-pollutant models, so the ordering is not a
potency comparison. The two Li et al.\ AMI estimates (28 Aug) are \emph{not} on a common
contrast: 1.019 is per IQR of total mass, 1.116 per quartile of a four-component mixture,
and the visual gap between them is partly a contrast artefact.}

\vspace{2mm}
\noindent\includegraphics[width=\linewidth]{\FIGDIR/f6_lifecourse.png}
\figcap{Pooled life-course coverage. Per-stage counts and annotations are summed and
merged across the seven daily issues rather than inherited from any one of them. The
working-age and 65+ windows carry this week --- CHARLS, the Korean screening cohort, the
Stockholm survey --- reversing last week's in-utero and early-childhood concentration,
consistent with \textit{Reproductive \& developmental} falling to 3 records.}

\vspace{2mm}
\noindent\includegraphics[width=\linewidth]{\FIGDIR/f4_heatmap.png}
\figcap{Subtopic $\times$ architecture over the pooled week. \textbf{The structural gap
this watch has reported every week is visible again}: the \textit{Sensing} and
\textit{Exposure} rows contain almost no cells that also carry a health endpoint, and the
cohort column contains almost no cells that measured their own exposure. Instrumentation
and epidemiology remain methodologically separate populations of paper. The
\textit{Metadata only} column, which spanned most subtopic rows last week, is now a thin
strip.}

\clearpage

% -------------------------------------------------------------- venues
\section*{\sffamily\bfseries\color{Deep}Venues, provenance and what the pipeline did}
\vspace{-1.5mm}{\color{Teal}\rule{\linewidth}{1.1pt}}\vspace{2.5mm}

\footnotesize
\textbf{Venues.} \textit{Atmos Chem Phys} and MDPI \textit{Toxics} led with 7 records
each, then \textit{Environ Pollut} and \textit{Environ Health (Wash)} with 5;
\textit{Zhongguo Fei Ai Za Zhi}, \textit{Atmosphere} and \textit{Front Public Health}
contributed 4 each. \textbf{Six preprints} entered the week (Research Square 3,
medRxiv 3) --- 5.5\% of the corpus. The \textit{Toxics} count is a single volume release
on 27 August rather than a sustained venue signal.

\vspace{1.2mm}
\textbf{Author clusters.} By surname the week is dominated by the expected Chinese
high-frequency names --- Chen (7 records), Zhang (6), Wang (4), Li (4), Yang (3), Wu (3)
--- but these are \emph{not} research groups: the corpus stores a short label, not an
affiliation or an ORCID, so surname frequency here measures nothing beyond the geography
already reported (China 27). \textbf{No genuine author-cluster analysis is possible with
the current record schema}, and this rollup does not attempt one. Adding a first-author
affiliation field is the cheapest fix and is logged as an open item.

\vspace{1.2mm}
\textbf{Two data defects were found by this rollup and are disclosed.} First,
\textbf{two records were published twice} --- article numbers
\texttt{2026.101.13} and \texttt{2026.102.16} under the
\texttt{10.3779/j.issn.1009-3419} prefix, both \textit{Zhongguo Fei Ai Za Zhi} ---
appearing in the 24 August issue and again on 25 August with differently-cased titles. The
DOI leg of \texttt{seen.json} should have caught this; it did not, which points at a
case-normalisation or write-ordering gap in the same-run screen. The pooled corpus is
therefore \textbf{109 records over 107 unique DOIs}, and every count in this issue is on
the 109 figure. Second, \textbf{five short-label collisions} were created within this
period --- two records each sharing \texttt{Chen et al.\ 2026e},
\texttt{Chen et al.\ 2026f}, \texttt{Li et al.\ 2026d}, \texttt{Wu et al.\ 2026c} and
\texttt{Zhang et al.\ 2026f}. Because the short label is the only identifier in the digest
and register, a collision makes two different papers indistinguishable. All five were
repaired before this rollup was built and the 28 and 29 August issues were recompiled and
re-shipped with the corrected labels. \textbf{Three older collisions remain in place}
(\texttt{Wang et al.\ 2026}, \texttt{Yang et al.\ 2026b}, and \texttt{Li \& Li 2026},
the last being the duplicate above) and are not retro-fixed, because rewriting a published
issue is worse than carrying a documented defect.

\vspace{1.2mm}
\textbf{Harvester state.} \textbf{OpenAlex returned HTTP 429 or 403 on every run of this
period} and should be treated as removed from the pipeline rather than degraded. Europe
PMC's \texttt{CREATION\_DATE} leg returned \textbf{zero} on both 28 and 29 August, which
is new --- it was returning archival back-scans in volume last week. arXiv contributed
nothing dated in window on any day of the period. Crossref by-ISSN carried the
instrumentation axis almost single-handedly: all three \textit{Atmos Chem Phys}
single-particle papers of 28 August, which PubMed structurally cannot index, arrived
through that leg alone. \textbf{Twenty-three design labels and two country mappings}
(\textit{Slovakia}, \textit{Italy}) were added to \texttt{plots.py} during the period, each
in the same edit that wrote them onto a record, per the 22 August guard; one further label
was found unmapped by this rollup and fixed here.

\vspace{1.2mm}
\textbf{Trial watch.} \texttt{state/trials.json} carries \textbf{23 registrations}, last
refreshed 22--26 August. Two entered in that window: \textbf{NCT05338242} (Monitoring
Environmental Exposures and Behavioral Change; interventional, completed, $n=30$) and
\textbf{NCT06247059} (Reducing Respiratory Virus Transmission in Bangladeshi Classrooms;
\textbf{withdrawn}, $n=0$). The withdrawal is the more informative of the two: classroom
filtration trials are exactly the interventional evidence this corpus keeps not
containing, and this one did not start. Of the 23 accumulated registrations, 21 are
interventional and \textbf{air cleaners or HEPA filtration are the intervention in 9 of
them} --- the registry pipeline is far more interventional than the published literature
this watch sees.

\vspace{3mm}
% -------------------------------------------------------------- gaps
\section*{\sffamily\bfseries\color{Deep}Open gaps}
\vspace{-1.5mm}{\color{Teal}\rule{\linewidth}{1.1pt}}\vspace{2.5mm}

\footnotesize
\begin{itemize}
\item \textbf{Nobody measured the exposure they analysed.} Of the 36 tier-A records this
week, every human-health one inherits a modelled exposure surface --- CMAQ, satellite
fusion, land-use regression or a reanalysis product --- and none reports that model's
error at the spatial scale of its own analysis. \textbf{Fifth consecutive week.} The one
partial exception cuts the other way: Wongsasuluk et al.\ (28 Aug) \emph{did} measure
their own exposure, with five low-cost DustBoy units in a road tunnel, and report no
co-location against a reference monitor --- so the risk estimates inherit an uncorrected
nephelometric bias in the highest-humidity, highest-loading regime there is.

\item \textbf{Constituent-resolved epidemiology still has no constituent-resolved
monitoring behind it.} The week's two strongest health results decompose PM$_{2.5}$ into
components taken from chemical transport or satellite products. Zhang et al.\ (28 Aug)
\emph{did} measure composition directly, at single-particle resolution, and has no health
endpoint. The two halves of the argument were published on the same day and do not cite
each other.

\item \textbf{Zero intervention studies in 109 records}, against 21 interventional
registrations sitting in the trial watch. The publication lag between the two is now the
single most quantifiable gap in this corpus.

\item \textbf{Ultrafine coverage doubled and is still negligible.} 7 of 109 records
(4 explicitly ultrafine, 3 size-resolved), up from 2 of 112. Best week on record, and
still under 7\%.

\item \textbf{Africa: 2 records in 109, both on the last day of the week.} The preceding
eight issues carried none. Both are useful --- a West Rand time-activity survey reporting
indoor fractions roughly 10 percentage points below USEPA defaults, and a meningitis-belt
dust review --- and neither involves an African-sited monitoring network.

\item \textbf{Endpoint canonicalisation, open since 15 August.} \texttt{ENDPOINT\_CANON}
still does not collapse synonymous endpoint strings, so \texttt{f2b}'s tail is inflated
and the ``no health endpoint'' count in this rollup (39) is assembled by rule in the
caption rather than read off the figure.
\end{itemize}

\clearpage

% -------------------------------------------------------------- digest
\section*{\sffamily\bfseries\color{Deep}Paper-level digest --- the week by subtopic}
\vspace{-1.5mm}{\color{Teal}\rule{\linewidth}{1.1pt}}\vspace{1mm}

{\fontsize{7.4}{9}\selectfont\color{Slate}Every record that appeared in a daily issue
keeps that issue's critical summary verbatim; nothing here is re-written or re-assessed.
Records surfaced without an abstract are shown as such. Ordered by subtopic, then by the
date of the issue that first carried the record.\par}

\input{digest_body}

\clearpage

% -------------------------------------------------------------- register
\section*{\sffamily\bfseries\color{Deep}Pooled corpus register}
\vspace{-1.5mm}{\color{Teal}\rule{\linewidth}{1.1pt}}\vspace{2.5mm}

{\fontsize{6.7}{8.4}\selectfont
\setlength{\tabcolsep}{2.4pt}
\begin{longtable}{@{}>{\RaggedRight}p{25mm}>{\RaggedRight}p{41mm}>{\RaggedRight}p{27mm}%
>{\RaggedRight}p{23mm}>{\RaggedRight}p{18mm}>{\RaggedRight}p{15mm}@{}}
\toprule
\textsf{\textbf{First author}} & \textsf{\textbf{Focus}} & \textsf{\textbf{Journal}} &
\textsf{\textbf{Design}} & \textsf{\textbf{Metric}} & \textsf{\textbf{Region}}\\
\midrule
\endfirsthead
\toprule
\textsf{\textbf{First author}} & \textsf{\textbf{Focus}} & \textsf{\textbf{Journal}} &
\textsf{\textbf{Design}} & \textsf{\textbf{Metric}} & \textsf{\textbf{Region}}\\
\midrule
\endhead
\multicolumn{6}{@{}l@{}}{\cellcolor{Mist}\textsf{\textbf{\textcolor{Deep}{Neuro / mental health}}}}\\[0.6mm]
\href{https://doi.org/10.1001/jamanetworkopen.2026.25015}{Dales} & Wildfire-sourced PM$_{2.5}$ and ED visits for migraine and other primary headache & \textit{JAMA Netw Open} & Case-crossover & PM$_{2.5}$ (wildfire) & Canada\\
\href{https://doi.org/10.1093/ageing/afag215}{Glover} & Umbrella meta-review of urban environmental exposures and cognitive health & \textit{Age Ageing} & Umbrella review / meta-analysis & PM$_{2.5}$ & Global\\
\href{https://doi.org/10.1016/j.nsa.2026.107019}{Goossens} & In utero air pollution exposure and pre-/postnatal brain development & \textit{Neurosci Appl} & Narrative review & PM$_{2.5}$ / PM$_{10}$ / PM$_{\mathrm{coarse}}$ & Global\\
\href{https://doi.org/10.1016/j.isci.2026.116617}{Ma} & PM$_{2.5}$-induced neurotoxicity and depression: exposure risks and mechanisms & \textit{iScience} & Narrative review & PM$_{2.5}$ + components & Global\\
\href{https://doi.org/10.3760/cma.j.cn112150-20250804-00752}{Pan} & Maternal PM exposure and offspring neurodevelopment: effects and mechanisms & \textit{Zhonghua Yu Fang Yi Xue Za Zhi} & Narrative review & PM (bulk) & China\\
\addlinespace[0.9mm]
\multicolumn{6}{@{}l@{}}{\cellcolor{Mist}\textsf{\textbf{\textcolor{Teal}{Cardiovascular \& metabolic}}}}\\[0.6mm]
\href{https://doi.org/10.3760/cma.j.cn112338-20251120-00843}{Huang} & PM$_{2.5}$ and dyslipidaemia in rural Chinese children; dietary effect modification & \textit{Zhonghua Liu Xing Bing Xue Za Zhi} & Cross-sectional (multistage) & PM$_{2.5}$ & China\\
\href{https://doi.org/10.32394/pe/218966}{Kazmierska} & Air pollution and cardiovascular disease: review of current data & \textit{Przegl Epidemiol} & Narrative review & PM$_{2.5}$ / PM$_{10}$ & Global\\
\href{https://doi.org/10.1007/s00125-026-06814-2}{Xi} & Air pollution, built environment and temperature vs. prediabetes phenotypes & \textit{Diabetologia} & Prospective cohort & PM$_{2.5}$ & Germany\\
\href{https://doi.org/10.1161/JAHA.125.046917}{Yang} & Short-term PM$_{2.5}$ and hospitalisation/ED use in atrial fibrillation & \textit{J Am Heart Assoc} & Case-crossover & PM$_{2.5}$ (wildfire-influenced) & USA\\
\addlinespace[0.9mm]
\multicolumn{6}{@{}l@{}}{\cellcolor{Mist}\textsf{\textbf{\textcolor{Sky}{Reproductive \& developmental}}}}\\[0.6mm]
\href{https://doi.org/10.1016/j.jes.2025.10.043}{Chen} & Prenatal air pollution, fetal growth trajectory phenotypes and child neurodevelopment & \textit{J Environ Sci (China)} & Prospective cohort & PM$_{2.5}$ / PM$_{10}$ & China\\
\href{https://doi.org/10.1016/j.ecoenv.2026.120566}{Huang} & Joint prenatal bisphenol + air-pollutant mixture exposure and preschool IQ & \textit{Ecotoxicol Environ Saf} & Prospective cohort (mixtures) & PM$_{2.5}$ / PM$_{10}$ & China\\
\href{https://doi.org/10.1007/s40572-026-00552-8}{Peterson} & Ambient air pollution, PAHs and pubertal development: critical review & \textit{Curr Environ Health Rep} & Critical review & PM$_{2.5}$ components & Global\\
\href{https://doi.org/10.1016/j.envres.2026.125313}{Xie} & Long-term PM$_{2.5}$ and timing of pubertal development in boys and girls & \textit{Environ Res} & Prospective cohort & PM$_{2.5}$ & China\\
\addlinespace[0.9mm]
\multicolumn{6}{@{}l@{}}{\cellcolor{Mist}\textsf{\textbf{\textcolor{Sage}{Respiratory \& allergic}}}}\\[0.6mm]
\href{https://doi.org/10.1007/s00484-026-03289-z}{Bayar Muluk} & Climate change and otolaryngology: airway disease impacts and emission reduction & \textit{Int J Biometeorol} & Narrative review & PM$_{2.5}$ & Global\\
\href{https://doi.org/10.4103/mgr.MEDGASRES-D-25-00164}{Liao} & Short-term air pollution, respiratory symptoms and lung function in asthma & \textit{Med Gas Res} & Panel study & PM$_{2.5}$ / PM$_{10}$ & China\\
\href{https://doi.org/10.3390/epidemiologia7040097}{Sanchez de Prada} & Ambient PM and asthma-coded admissions in children 0-3 y, Spain & \textit{Epidemiologia} & Ecological panel & PM$_{2.5}$ / PM$_{10}$ & Spain\\
\href{https://doi.org/10.1016/j.jhazmat.2026.143076}{Yang} & Allergenic vs. toxic risk of airborne seed hairs compared with ambient PM$_{2.5}$ & \textit{J Hazard Mater} & Field measurement & PM$_{2.5}$ + bioaerosol & China\\
\addlinespace[0.9mm]
\multicolumn{6}{@{}l@{}}{\cellcolor{Mist}\textsf{\textbf{\textcolor{Violet}{Mechanistic toxicology}}}}\\[0.6mm]
\href{https://doi.org/10.1016/j.jes.2025.12.074}{Ding} & PM$_{2.5}$ promotes vascular calcification via miR-27b-5p/HES1 signalling & \textit{J Environ Sci (China)} & In vitro + organoid + murine & PM$_{2.5}$ & China\\
\href{https://doi.org/10.1007/s00360-026-01698-5}{Maraschi} & Osmoregulatory disruption and metal bioaccumulation from settleable atmospheric PM & \textit{J Comp Physiol B} & Controlled exposure (ecotox) & Settleable APM & Brazil\\
\href{https://doi.org/10.1016/j.jes.2025.11.018}{Wu} & Macrophage NLRP3 inflammasome in real-ambient PM-induced hepatic glucose dysregulation & \textit{J Environ Sci (China)} & Real-ambient murine + co-culture & Real-ambient PM & China\\
\addlinespace[0.9mm]
\multicolumn{6}{@{}l@{}}{\cellcolor{Mist}\textsf{\textbf{\textcolor{Clay}{Occupational \& indoor}}}}\\[0.6mm]
\href{https://doi.org/10.3389/fmedt.2026.1796828}{Chen} & Nested suction instrument to cut intraoperative surgical-smoke PM$_{2.5}$ & \textit{Front Med Technol} & Bench / device experiment & PM$_{2.5}$ (surgical smoke) & China\\
\href{https://doi.org/10.1093/annweh/wxag055}{Cole} & Determinants of respirable crystalline silica exposure in Australian tunnelling & \textit{Ann Work Expo Health} & Observational exposure study & Respirable dust / RCS & Australia\\
\href{https://doi.org/10.1016/j.jhazmat.2026.143081}{Ding} & Stage-resolved hair Zn-isotope dynamics as a firefighting combustion-exposure tracer & \textit{J Hazard Mater} & Repeated-measures biomarker & Combustion-derived PM & China\\
\href{https://doi.org/10.3390/jox16040119}{Ramsden} & Xenobiotic hazards in aircraft cabin air & \textit{J Xenobiot} & Review + risk economics & Ultrafine particles & Global\\
\addlinespace[0.9mm]
\multicolumn{6}{@{}l@{}}{\cellcolor{Mist}\textsf{\textbf{\textcolor{Amber}{Sensing, forecasting \& instrumentation}}}}\\[0.6mm]
\href{https://doi.org/10.1016/j.jenvman.2026.130515}{Shi} & Adaptive spatiotemporal graph transformer for multi-site multi-step PM$_{2.5}$ forecasting & \textit{J Environ Manage} & Deep-learning model & PM$_{2.5}$ & China / India\\
\href{https://doi.org/10.3390/bios16070380}{Tsuchiya} & Seasonal and indoor-outdoor characteristics of fluorescent aerosol particles in offices & \textit{Biosensors} & Real-time sensor campaign & Bioaerosol / coarse & Japan\\
\addlinespace[0.9mm]
\multicolumn{6}{@{}l@{}}{\cellcolor{Mist}\textsf{\textbf{\textcolor{Amber}{Exposure assessment \& modelling}}}}\\[0.6mm]
\href{https://doi.org/10.1029/2025GH001621}{Solorzano-Araque} & Multiscale exposure modelling: LOTOS-EUROS + mobility surveys + particle toxicity & \textit{GeoHealth} & CTM + mobility integration & PM$_{2.5}$ / PM$_{10}$ & Colombia\\
\addlinespace[0.9mm]
\multicolumn{6}{@{}l@{}}{\cellcolor{Mist}\textsf{\textbf{\textcolor{Coral}{Burden, policy \& mitigation}}}}\\[0.6mm]
\href{https://doi.org/10.1002/hsr2.72886}{Karamzad} & Ischaemic heart disease burden in MENA 1990-2021; PM attributable fraction 28.2\% & \textit{Health Sci Rep} & GBD secondary analysis & Ambient PM & MENA\\
\href{https://doi.org/10.1371/journal.pone.0347989}{Rees} & HEPA filtration to reduce respiratory infection in care homes & \textit{PLoS One} & RCT process evaluation & Filtered PM & UK\\
\href{https://doi.org/10.1016/j.jes.2026.01.009}{Shan} & Policy-driven emission reductions and post-pandemic rebound, high-resolution inventories & \textit{J Environ Sci (China)} & Emission inventory & PM emissions & China\\
\addlinespace[0.9mm]
\multicolumn{6}{@{}l@{}}{\cellcolor{Mist}\textsf{\textbf{\textcolor{Slate}{Other clinical endpoints}}}}\\[0.6mm]
\href{https://doi.org/10.3390/dj14070440}{Lupita} & PM$_{2.5}$, preventive-care capacity and caries/edentulism burden across the EU & \textit{Dent J} & Ecological panel & PM$_{2.5}$ & EU-27\\
\href{https://doi.org/10.1016/j.jdin.2026.06.008}{Park} & Long-term PM exposure and risk of androgenetic alopecia & \textit{JAAD Int} & Nationwide cohort & PM$_{2.5}$ / PM$_{10}$ & South Korea\\
\href{https://doi.org/10.3171/2026.3.PEDS2615}{Sun} & Environmental exposure and area deprivation vs. paediatric hydrocephalus outcomes & \textit{J Neurosurg Pediatr} & Registry cohort + geospatial & PM$_{2.5}$ (satellite-derived) & USA\\
\addlinespace[0.9mm]
\bottomrule
\end{longtable}}


\vspace{4mm}
\begin{keybox}
{\sffamily\bfseries\footnotesize\color{Deep}Provenance and method}\\[1.2mm]
{\fontsize{7.2}{8.8}\selectfont
Period \textbf{2026-08-23 to 2026-08-29}, seven daily issues pooled.
\textbf{This rollup queried no API.} It aggregates \texttt{state/corpus/*.json},
\texttt{state/metrics.csv} and \texttt{state/trials.json} over
\texttt{PMRW\_START}--\texttt{PMRW\_END}; every record was screened, summarised and
DOI-gated at daily-issue time and is reproduced verbatim. \textbf{109 records, 107 unique
DOIs} --- the two duplicates are named on the venues page. Daily totals: 9, 9, 14, 14, 36,
19, 8. \textbf{Week-over-week comparisons} are against the 16--22 August rollup
(112 records) and, where a longer run is cited, the 8 and 15 August rollups.
\textbf{Effect estimates} are pooled as recorded, on heterogeneous contrasts, with no
meta-analysis. \textbf{Figures} are regenerated over the pooled range, not composited from
the daily images. \textbf{Labels} --- subtopic, design, particle-metric, geography --- are
single-label by dominant contribution, assigned at daily-issue time, and are a judgement
rather than a controlled vocabulary.\par}
\end{keybox}

\end{document}
